\begin{table}[htbp]
\centering
\caption{Definitions and clinical interpretation of study variables.}
\label{tab:variable_definitions}
\small
\begin{tabular}{p{3.2cm}p{2.1cm}p{8.6cm}}
\toprule
Variable & Type & Description \\
\midrule
age & Continuous & Age of the patient in years; a general demographic indicator associated with frailty, comorbidity burden, and prognosis. \\
anaemia & Binary & Presence of anaemia (0 = no, 1 = yes); reflects reduced oxygen-carrying capacity and may indicate worse overall clinical status. \\
creatinine\_phosphokinase & Continuous & Serum creatinine phosphokinase level (mcg/L); a laboratory marker that may reflect muscle injury or systemic stress. \\
diabetes & Binary & History of diabetes mellitus (0 = no, 1 = yes); an important cardiometabolic comorbidity relevant to heart failure progression. \\
ejection\_fraction & Continuous & Left ventricular ejection fraction (\%); a central measure of cardiac systolic function and heart failure severity. \\
high\_blood\_pressure & Binary & Presence of hypertension (0 = no, 1 = yes); a common cardiovascular comorbidity and contributor to cardiac remodeling. \\
platelets & Continuous & Platelet count (kiloplatelets/mL); a hematologic variable that may reflect inflammatory, thrombotic, or systemic disease processes. \\
serum\_creatinine & Continuous & Serum creatinine concentration (mg/dL); an indicator of renal function, which is closely linked to prognosis in heart failure. \\
serum\_sodium & Continuous & Serum sodium concentration (mEq/L); a biochemical marker relevant to volume status and neurohormonal dysregulation. \\
sex & Binary & Biological sex (0 = female, 1 = male); included as a demographic characteristic that may relate to disease presentation and outcome. \\
smoking & Binary & Smoking status (0 = no, 1 = yes); a behavioral cardiovascular risk factor that may influence long-term clinical status. \\
time & Continuous & Follow-up time in days; an observation-related variable that is highly informative for outcome discrimination but should be interpreted cautiously in etiologic risk-factor discussion because it reflects follow-up structure as well as event occurrence. \\
\texttt{DEATH\_EVENT} & Binary outcome & Study endpoint (0 = survivor, 1 = death) used as the target variable for prediction and comparative factor analysis. \\
\bottomrule
\end{tabular}
\end{table}
